% Theme: Archive  --  best for history, famous equations, quotes, timelines
% Build:  ./build.sh example-archive.tex   (or compile on Overleaf)
\documentclass[12pt]{beamer}
\usetheme{igarchive}

\ighandle{@quantum_danesh}
\igseries{Historia cuántica}
\igissue{02}

\begin{document}

\coverslide{De los archivos}
  {1935, el año del entrelazamiento}
  {De una objeción de Einstein a un premio Nobel.}

\begin{frame}{La objeción de EPR}{Mayo de 1935}
\lettrine{E}{instein}, Podolsky y Rosen publican en \emph{Physical Review}
una pregunta incómoda: ¿puede considerarse \hl{completa} la descripción
cuántica de la realidad?

Su ejemplo: dos partículas tan correlacionadas que, al medir una, conoces
la otra sin tocarla. Para ellos, eso probaba que a la teoría le faltaba
algo.
\vfill
{\footnotesize\color{igMuted}Phys.\ Rev.\ \textbf{47}, 777 (1935).\par}
\end{frame}

\begin{frame}{Schrödinger le pone nombre}{1935}
Ese mismo año, Schrödinger llama al fenómeno \key{Verschränkung}, en
inglés \key{entanglement}, y lo describe así:
\begin{pullquote}{E.~Schrödinger, 1935 (traducción)}
Yo no lo llamaría \emph{un} rasgo, sino \emph{el} rasgo característico de
la mecánica cuántica, el que impone su ruptura total con las líneas del
pensamiento clásico.
\end{pullquote}
\end{frame}

\begin{frame}{Un estado inseparable}{Hoy}
\[
  \ket{\Phi^+} = \frac{1}{\sqrt2}\big(\ket{00} + \ket{11}\big)
\]
Si fuera un producto
\[
  (\alpha_0\ket0+\alpha_1\ket1)\otimes(\beta_0\ket0+\beta_1\ket1)
\]
haría falta
\[
  \alpha_0\beta_0 = \alpha_1\beta_1 = \tfrac{1}{\sqrt2},
  \qquad
  \mhl{\alpha_0\beta_1 = \alpha_1\beta_0 = 0}.
\]
Imposible: la segunda anula un factor de la primera. Está \hl{entrelazado}.
\end{frame}

\begin{frame}{Del papel al laboratorio}{1935 -- 2022}
\begin{center}
\begin{tikzpicture}[x=1mm,y=1mm]
  \draw[igFG, line width=0.5pt] (0,2) -- (0,-52);
  \foreach \y/\yr/\txt in {0/1935/{EPR y Schrödinger: el problema},
      -10/1964/{Bell: una desigualdad medible},
      -20/1972/{Freedman y Clauser: primer test},
      -30/1982/{Aspect y su equipo, en Orsay},
      -40/2015/{Primeros tests sin lagunas},
      -50/2022/{Nobel: Aspect, Clauser, Zeilinger}}{
    \fill[igBG] (0,\y) circle (1.5);
    \draw[igAccent, line width=0.6pt] (0,\y) circle (1.5);
    \fill[igAccent] (0,\y) circle (0.6);
    \node[anchor=base west, text=igAccent, font=\large\itshape] at (3.5,\y-1.3) {\yr};
    \node[anchor=base west, text=igFG, font=\small] at (16.5,\y-1.3) {\txt};
  }
\end{tikzpicture}
\end{center}
\end{frame}

\begin{frame}{La desigualdad CHSH}{1969}
Con dos medidas posibles a cada lado,
\[
  S = \expect{A_0B_0} + \expect{A_0B_1} + \expect{A_1B_0} - \expect{A_1B_1}.
\]
Cualquier teoría local y realista cumple $|S| \le 2$. La mecánica cuántica
llega más lejos:
\keyeq{|S|_{\max} = \mhl{2\sqrt{2}}}
Es la cota de Tsirelson (1980), y los experimentos le dan la razón.
\end{frame}

\closingslide{1935 sigue abierto}
  {Próximo capítulo: 1982, el experimento de Aspect paso a paso.}

\end{document}
